\section{生产实习概况}

\subsection{实习日程安排}
本次实习分两周进行，第一周在南京浦镇，主要学习轨道车辆，尤其是城轨车辆的制造工艺，并在工厂内进行观摩学习；第二周在上海各轨交基地进行，主要学习的是地铁列车、国铁动车组的运维知识。\par{}
具体日程安排根据实习安排表和实际情况做了修正，如表~\ref{tab:schedule}所示：\par{}
% This table intentionally uses \hline and \cline to preserve merged-cell borders.
% chktex-file 1
% chktex-file 8
% chktex-file 44

\begingroup
\hbadness=10000
\setlength{\tabcolsep}{2pt}
\setlength{\extrarowheight}{1pt}
\renewcommand{\arraystretch}{1.10}
\centering
\reportSong\fontsize{8pt}{10pt}\selectfont
\newcommand{\compactcell}[1]{\begingroup\renewcommand{\arraystretch}{0.78}\begin{tabular}[c]{@{}c@{}}#1\end{tabular}\endgroup}
\newcommand{\tallcell}[1]{\parbox[c][11.5mm][c]{\linewidth}{\centering #1}}
\newcommand{\scheduleDate}[2]{\Block[c,m]{#1-1}{#2}}
\newcommand{\scheduleUnit}{\Block[c,m]{10-1}{\compactcell{中车南京浦镇车辆\\有限公司}}}
\begin{table}[H]
  \centering
  \vspace{-2mm}
  \caption{实习日程安排}\label{tab:schedule}
  \vspace{-3mm}
  \reportSong\fontsize{8pt}{10pt}\selectfont
  \makebox[\linewidth][c]{
    \begin{NiceTabular}{|>{\centering\arraybackslash}m{26mm}|>{\centering\arraybackslash}m{16mm}|>{\centering\arraybackslash}m{14mm}|>{\centering\arraybackslash}m{72mm}|>{\centering\arraybackslash}m{23mm}|}
      \hline
      \multicolumn{1}{|c|}{\reportHei\fontsize{9pt}{11pt}\selectfont 实习单位} &
      \multicolumn{2}{c|}{\reportHei\fontsize{9pt}{11pt}\selectfont 时间}    &
      \multicolumn{1}{c|}{\reportHei\fontsize{9pt}{11pt}\selectfont 教学内容}  &
      \multicolumn{1}{c|}{\reportHei\fontsize{9pt}{11pt}\selectfont 地点}                                                                                                             \\
      \hline
      \scheduleUnit                                                        & \scheduleDate{4}{7月6日}  & \Block[c,m]{2-1}{上午} &
      浦镇公司产品创新、升级与研发体系介绍                                                   &
      \Block[c,m]{4-1}{浦镇文体报告厅}                                                                                                                                                     \\
      \cline{4-4}
                                                                           &                         &                      & 城轨车体概述及车体工艺制造技术介绍          &                          \\
      \cline{3-4}
                                                                           &                         & \Block[c,m]{2-1}{下午} & 转向架工艺设计流程介绍和焊接、组装智能产线介绍    &                          \\
      \cline{4-4}
                                                                           &                         &                      & 城轨车辆总装生产工艺技术概述             &                          \\
      \cline{2-5}
                                                                           & \scheduleDate{2}{7月7日}  & 上午                   & 企业文化展厅参观                   & 浦镇企业展厅                   \\
      \cline{3-5}
                                                                           &                         & 全天                   & 转向架分厂现场实习                  & 浦镇转向架分厂                  \\
      \cline{2-5}
                                                                           & 7月8日                    & 全天                   & 车体分厂现场实习                   & 浦镇车体分厂                   \\
      \cline{2-5}
                                                                           & 7月9日                    & 全天                   & \Block[c,m]{2-1}{总装分厂现场实习} & \Block[c,m]{2-1}{浦镇总装分厂} \\
      \cline{2-3}
                                                                           & \scheduleDate{2}{7月10日} & 上午                   &                            &                          \\
      \cline{3-5}
                                                                           &                         & 下午                   & \centering 交流活动            & 浦镇文体报告厅                  \\
      \hline
      \tallcell{\compactcell{上海中车申通轨道                                                                                                                                               \\交通车辆有限公司}} & \tallcell{7月13日}                   & \tallcell{上午}                   & \tallcell{参观赛车场车辆段维修中心，了解列车检修流程}      & \tallcell{\compactcell{上海地铁赛车场\\车辆段}} \\
      \cline{1-5}
      \tallcell{\compactcell{上海地铁维护保障                                                                                                                                               \\有限公司车辆分公司}} & \tallcell{7月14日}                   & \tallcell{上午}                   & \tallcell{参观梅陇维保中心，学习数字化运维与车辆检测实验室}   & \tallcell{\compactcell{上海地铁梅陇\\维保中心}} \\
      \cline{1-5}
      \tallcell{\compactcell{上海阿尔斯通                                                                                                                                                 \\交通设备有限公司}}     & \tallcell{7月15日}                   & \tallcell{上午}                   & \tallcell{听取企业与车辆设计讲座，了解新线市域车设计}      & \tallcell{\compactcell{上海轨道交通\\设备发展有限公司}} \\
      \cline{1-5}
      \tallcell{\compactcell{中国铁路上海局集团                                                                                                                                              \\有限公司上海动车段}} & \tallcell{7月16日}                   & \tallcell{上午}                   & \tallcell{参观动车组高级修场，学习三级修与部件检修流程}     & \tallcell{\compactcell{上海动车段\\高级修场}} \\
      \hline
    \end{NiceTabular}
  }
\end{table}
\endgroup


\subsection{实习单位概况}
本次实习前往了多个轨道车辆设计、制造、运维、管理单位，下面是各单位基本信息；由于保密要求，后文中涉及某些单位的内容无任何现场照片，相关部分的图像资料均来源于一手记录或网络资料。

\subsubsection{中车南京浦镇车辆有限公司}
中车南京浦镇车辆有限公司位于南京市浦口区泰山街道浦珠北路68号，分为老厂区和新厂区两部分，其核心业务就是轨道交通装备制造业务，支柱业务是运营、维保、修理等延伸业务，支撑业务是智能制造、智能产品、智能运营等新产业业务。目前，地铁、市域车设计、制造、总装和调试等环节主要由新厂区承担，而老厂区主要承担国铁客车车厢的修缮工作。\par{}
该公司具备较完整的研发能力，包括车体、转向架、TCMS等关键系统与产品的研发，网络软件的研发，静强度、碰撞、动应力、模态等的多学科仿真设计，系统工程设计，工业设计，基础研究及试验等能力，涵盖车辆设计、生产与制造全流程的研发任务。同时其产品平台也十分多样，如系列化动力集中/分散型动车组、国铁客车、城轨车辆及各类工程车产品等。\par{}
人才培养方面，新员工经过一年左右的磨合后可选择两条发展通道，分别是行政管理与专业技能通道；招聘技术人员以校招为主，共分为L1\textasciitilde{}L10等级，校招技术人员一般为L8或L9级，后面每2或3年进行评级。\par{}

\subsubsection{上海中车申通轨道交通车辆有限公司（赛车场车辆段）}
上海中车申通轨道交通车辆有限公司下属赛车场车辆段位于上海市嘉定区安亭镇伊宁路2620号，由11号线上海赛车场站花桥方向交叉渡线两端引出，主要承担上海地铁11号线的修缮维护工作，由总体拆卸、车体设备维护、转向架拆卸与维护、轮对拆卸与维护、电机维护、各部件到总体依次组装的过程构成。\par{}
其在配线图中的位置与现场外景如图\ref{fig:saichang-maintenance-base}\sourcecite{1}所示。\par{}
\reportFigure{8}

\subsubsection{上海地铁维护保障有限公司车辆分公司（梅陇基地）}
\noindent
\begin{minipage}[t]{0.48\textwidth}
  \vspace{-3mm}
  \qquad{}上海地铁维护保障有限公司车辆分公司下属梅陇基地位于上海市徐汇区老沪闵路1号，介于上海地铁1号线地面段与金山铁路之间，如图\ref{fig:meilong-base}\sourcecite{1}所示，可视为上海地铁1号线的里程起点，主要承担上海轨道交通车辆运营管理、维护保障、设备试验等相关工作。\par{}
  \qquad{}面对上海日益发展完善的轨道交通，在车辆业务调整、路网规模扩大、检修模式变化的考量之下，该公司建立了\par{}
\end{minipage}
\hfill
\begin{minipage}[t]{0.48\textwidth}
  \reportInlineFigure{9}
\end{minipage}
\par\medskip\noindent
1~+~6~+~33~+~X的指挥布局，即1个车辆数字化运维管理中心、6个车辆数字化运维支持中心、33个车辆控制中心、多个数字化运营协调中心的布局，有助于提高车辆运维组织效率。\par{}
其人才培养体系将人才分为5类对象：新入职员工、青年人才、管理人员、生产人员、拔尖人才，有管理序列、专业技术序列、生产序列3种培养路径，将员工的发展分为职业导航期、轮岗锻炼期和专业进阶期这3个阶段，并设置11个等级。

\subsubsection{上海轨道交通设备发展有限公司（阿尔斯通）}
上海轨道交通设备发展有限公司位于上海市闵行区东川路3999号闵行开发区站西侧，该公司由中车长客、上海电气等共同持股，融合阿尔斯通等品牌、资源、技术优势，下属上海中车长客交通设备有限公司（原SATCO）、上海阿尔斯通交通电气有限公司（SATEE）、工程公司等子公司、联营公司或独立分公司，业务涉及车辆设计、关键系统设备、升级改造及工程服务等。\par{}
原SATCO具有从车辆设计，到组装、涂装、总装、调试、质保、维修的一系列完整产业链，参与了多种上海地铁车型的设计；而SATEE则负责生产牵引系统设备、辅助变流器及其相关电气设备，是国内具有投标资质的定点生产牵引传动与控制系统的企业；而工程公司则负责设计制造安装屏蔽门、安全门等产品，还具备供电、车站、车辆段等设备集成与施工安装的相关能力。\par{}
该单位位于闵行开发区站西侧，其在配线图中的位置如图\ref{fig:satee-position}\sourcecite{1}所示。\par{}
\reportFigure{10}

\subsubsection{中国铁路上海局集团有限公司上海动车段（高级修场）}
中国铁路上海局集团有限公司上海动车段位于上海市嘉定区江桥镇老翔黄公路600号，承担动车组运用检修和高级修相关任务，本次实习参观的是动车组高级修场，位于南翔编组站南侧，是全路率先建成的四个高修场地之一；其通过西走行线与上海虹桥动车所、上海虹桥站连接，通过东走行线与南翔动车所、上海西站、上海站连接。\par{}
上海动车段动车组高级修场在配线图中的位置如图\ref{fig:shanghai-emu-depot}\sourcecite{1}所示。\par{}
\reportFigure{11}


